\documentclass[11pt, a4paper]{report}

\usepackage[utf8]{inputenc}
\usepackage[T1]{fontenc}
\usepackage[french]{babel}

% Appel des packages dans le dossier preambule
\def\packagepath{../../../../preambule}   % path du package princial

\usepackage{\packagepath/page_layout} % <--- Nouveau
\usepackage{\packagepath/config_info}
\usepackage{\packagepath/cadres_cours}
\usepackage{\packagepath/zones_reponse}
\usepackage{\packagepath/config_ds}
\usepackage{\packagepath/config_maths}

% --- PILOTAGE ---
\def\visibleornot{visible}    % visible or invisible


\def\documentpath{./Sources_Latex}
\graphicspath{{./Images/}}


\def\ptsexoA{0}

\begin{document}

%%%%%%%%%%%%%%%%%%%%%%%%%%%
\def\level{BTS}  % Classe à droite
\def\course{CIEL 1}
\def\chaptername{Nombes complexes} % Titre au centre
\def\docname{C112~--~TP1}        % Identifiant à gauche

\pagestyle{DS_LP}
%%%%%%%%%%%%%%%%%%%%%%%%%%%%%%%%%

\NomPrenomNote{}

%%=============================================================

\section*{PARTIE A\@: Création et manipulation de base}
\medskip




\begin{enumerate}
    \item Dans la barre de saisie, tapez: \verb|z_1 = 3 + 2|.
    \item Déplacez le point $z_1$ à la souris et observez comment sa valeur change dans la fenêtre d'algèbre.
    \item Créez son conjugué en tapant: \verb|z_2 = conjugate(z_1)|.
    
    Quel type de symétrie observe-t-on entre $z_1$ et $z_2$?

    \begin{blocvisible}[height=2]
        Le point $z_2$ est le symétrique de $z_1$ par rapport à l'axe des abscisses (l'axe réel). En effet, si $z_1 = a + \i b$, alors son conjugué $z_2 = a - \i b$, ce qui correspond à une réflexion par rapport à l'axe réel dans le plan complexe.
    \end{blocvisible}    
    \item Par le calcul déterminer le module de $z_1$ et de $z_2$. 
    \begin{blocvisible}[height=4]
        Le module d'un nombre complexe $z = a + \i b$ est donné par la formule $|z| = \sqrt{a^2 + b^2}$. 

        Pour $z_1 = 3 + 2\i$, on a:
        \[|z_1| = \sqrt{3^2 + 2^2} = \sqrt{9 + 4} = \sqrt{13}\]

        Pour $z_2 = 3 - 2\i$, on a:
        \[|z_2| = \sqrt{3^2 + {(-2)}^2} = \sqrt{9 + 4} = \sqrt{13}\]

        Ainsi, les modules de $z_1$ et $z_2$ sont égaux: $|z_1| = |z_2| = \sqrt{13}$.
    \end{blocvisible}
    \item Vérifier la valeur du module de $z_1$ en tapant: \verb|r = abs(z_1)|. Faites la même chose pour $z_2$.

\end{enumerate}

\section*{Partie B\@: Forme Trigonométrique dynamique}


\begin{enumerate}[resume]
    \item Créez un curseur $a$ (theta) allant de $0$ à $2\pi$ (en radians).
    \item Créez un curseur \texttt{mod} (le module) allant de $0$ à $10$.
    \item Saisissez le point $z_3$ sous sa forme trigonométrique: \verb|z_3 = mod * (cos(a) + i * sin(a))|.
    \item  Fixez \texttt{mod = 4} et faites varier \verb|a|. 
    \item Activez la \og{}Trace\fg{} sur $z_3$. Quelle figure géométrique obtenez-vous?
\begin{blocvisible}[height=2]
    En faisant varier l'argument $a$ de $0$ à $2\pi$ tout en maintenant le module constant à $4$, le point $z_3$ décrit un cercle de rayon $4$ centré à l'origine du plan complexe. Ainsi, la figure géométrique obtenue est un cercle de rayon $4$.

\end{blocvisible}     
\end{enumerate}

\section*{Partie C\@: Somme de vecteurs complexes}

\begin{enumerate}[resume]
    \item Créez deux points libres $z_A$ et $z_B$.
    \item Saisissez \verb|Somme = z_A + z_B|.
    \item Tracez les vecteurs: \verb|u = Vector(z_A)| et \verb|v = Vector(z_B) et w = Vector(Somme)|.
    \item Faites varier les points $z_A$ et $z_B$. Que constatez-vous?

    Vérifiez que la position du point \verb|Somme| correspond bien à la règle du parallélogramme apprise en cours.
    \begin{blocvisible}[height=2]
    Le vecteur \verb|w| représente la somme des vecteurs \verb|u| et \verb|v|. En faisant varier les points $z_A$ et $z_B$, on observe que le vecteur \texttt{w} correspond à la diagonale du parallélogramme formé par les vecteurs \texttt{u} et \texttt{v}, ce qui illustre la règle de la somme des vecteurs dans le plan complexe.
\end{blocvisible}
\end{enumerate}
    \end{document}







